\chapter{作业押题-成长感悟类}

\begin{itemize}
    \item 《那一刻，我读懂了\_\_》
    \item 《这样的挫折让我成长》
    \item 《藏在\_\_里的蜕变》
    \item 《为了心中的那束光》
\end{itemize}


\section{那一刻，我懂得了坚持}

人生之路，犹如一幅波澜壮阔的画卷，上面布满了无数的挑战与艰难险阻。在面对这些困境之时，我们常常会陷入迷茫与无助的境地，甚至萌生出放弃的念头。然而，总有那么一个独特而难忘的时刻，如同夜空中最璀璨的星辰，照亮我们的心灵，让我们深刻地领悟到坚持的强大力量。那一刻，我真正懂得了坚持的含义。

那是一次学校精心组织的长跑比赛。我，一个向来对运动不太擅长的人，原本对这场比赛并未抱有过多的期望。然而，在同学们热情洋溢的鼓励下，我最终决定勇敢地参加这次比赛，挑战一下自我。

比赛的那一天，阳光灿烂得如同金子般洒落在大地上，微风轻柔地拂过脸庞。操场上，同学们个个摩拳擦掌，为比赛做着最后的准备。我紧张而又忐忑地站在起跑线上，心脏如同擂鼓一般“咚咚”直响。随着一声清脆的发令枪响，同学们犹如离弦之箭般飞速冲了出去。我也不甘落后，紧紧地跟在队伍的后面。

可没过多久，我就开始感到体力渐渐不支。我的呼吸变得急促而紊乱，仿佛拉风箱一般；双腿也仿佛灌了铅似的，沉重得难以抬起。看着前面的同学越跑越远，我的心中涌起一股强烈的想要放弃的冲动。但是，就在这时，我想起了自己报名参加比赛的初衷，想起了同学们那充满期待的眼神和鼓励的话语。我暗暗告诉自己，绝对不能就这样轻易放弃，一定要咬牙坚持下去。

于是，我努力调整自己的呼吸，放慢脚步，一步一步艰难地向前跑去。每迈出一步，我都能清晰地感受到自己的身体在承受着巨大的压力，仿佛随时都会崩溃。然而，我始终没有停下自己的脚步。我深知，只要我坚持下去，就一定能够完成这场比赛。

在那漫长而又似乎没有尽头的跑道上，我不断地给自己加油打气，不断地挑战着自己身体的极限。终于，我看到了那条通往终点的线。那一刻，我的心中充满了难以言表的喜悦和强烈的成就感。我用尽最后一丝力气，奋力冲过了终点线。

那一刻，我懂得了坚持。坚持，就是在面对重重困难和挫折时，绝不放弃、绝不退缩，勇敢而坚定地向前迈进。坚持，让我战胜了自己内心的软弱和恐惧，让我完成了原本以为自己根本无法完成的任务。从那以后，每当我在生活中遇到困难想要放弃的时候，我都会不由自主地想起那一刻，想起自己在长跑比赛中坚持到底的情景。那一刻，已然成为我人生中最为宝贵的财富，它将永远激励着我在人生的道路上勇敢无畏地前行。

\section{我终于听懂了清晨}

以前，我讨厌清晨。

尤其是初三的清晨。天还没完全亮，闹钟已经响了。那声音像是从另一个世界砸下来的，硬生生把我从梦里拽出来。窗外的鸟叫，在我听来不是欢快，而是吵闹——它们凭什么那么高兴？妈妈煮粥的声音从厨房传来，锅盖碰撞的轻响，也不像温暖，更像一种提醒：你又该去受苦了。

那段时间，我总觉得生活被课表和倒计时填满了。每天背单词、刷题、改错，一遍一遍，像在一条看不见尽头的隧道里奔跑，跑得气喘吁吁，却不知道尽头在哪里，也不知道为什么要跑。

一次早读，我因为前一晚睡得太晚，整个人昏昏沉沉，老师点我起来回答问题，我一个字也说不出。全班安静了两秒，那两秒比两小时还长。我心里又羞又烦，放学回家，终于忍不住向妈妈爆发：“为什么每天都这么累？努力到底有什么用？凭什么别人可以睡到七点，我六点就要爬起来？”

妈妈没有立刻回答。她看了我一会儿，只说了一句：“明天早上，我带你下楼买个豆浆。”

第二天清晨，她果然提前叫醒了我。我不情愿地跟着她下楼，嘴里还嘟囔着“困死了”。小区还很安静，路灯没有完全熄灭，橘黄色的光晕落在潮湿的地面上。天边却已经泛出淡淡的蓝，像一块被水洗过的绸布，边缘透出一点点金。

早餐摊的阿姨正在支炉子，白色的蒸汽从锅盖边缘冒出来，热腾腾地升向微凉的天空。保安叔叔拿着大扫帚，一下一下清理昨夜的落叶，扫帚划过水泥地面的声音，沙沙的，像在跟谁低声说话。几个背书包的学生匆匆走过，嘴里还念着英语单词，声音含糊却执着，像在跟时间赛跑。

妈妈指着这些人，语气很轻：“你看，清晨不是只有你一个人在努力。”

我愣住了。

原来，在我抱怨疲惫的时候，这座城市早已醒来。有人为了生计早起，有人为了责任坚守，有人为了梦想奔跑。他们谁也没有站在山顶上，谁也没有被多少人看见，但他们都在自己的轨道上，认真活着。

我们买了两杯热豆浆。纸杯捧在手里，温度从掌心一点点渗进来，沿着手臂慢慢往上走。回去的路上，一缕阳光越过楼顶，斜斜地落在湿润的地面上，把水洼照得亮晶晶的。我忽然听见鸟叫——那些我曾经嫌吵的声音，此刻变得清亮而又有层次，一声高，一声低，像在为刚刚开始的这一天试音。

那一刻，我第一次觉得，清晨并不冷漠。它只是把希望藏得很浅，浅到只要你肯抬头，就一定能看见。

那天之后，我把闹钟铃声换成了一段轻柔的鸟鸣。它响起的时候，我仍会眯着眼赖上几秒，但心里不再只剩抗拒。因为我知道，在这座城市的很多扇窗户后面，有无数和我一样的人，正在同一时刻醒来。我们互不相识，却共享同一个清晨。

后来，我又去买了许多次豆浆。每一次，我都会多看几眼那些在清晨里忙碌的人。环卫工扫过的地方，路面干净得像被洗过；早餐摊的蒸汽升起来，在朝阳里变成金色的雾；那些背着书包念念有词的学生，从一个变成两个，从两个变成一群，像溪流汇入河流，慢慢填满整条街道。

我终于听懂了清晨。

它说，努力也许不会立刻开花，也许你永远站不到山顶，也许你的光芒永远只是微光，但那又怎样？每一次准时醒来，每一次咬牙坚持，都是在为这个世界点亮一盏小小的灯。灯多了，天就亮了。

疲倦、误解、昨夜的沮丧，都可以在第一缕光里慢慢放下。只要心中还有清晨，前路就不会真正黑暗。

并非站在山顶，才能被看见。那些在山脚下、在街角、在路灯还没有熄灭时就醒来的人，他们散发的微光，同样照亮了人间。而我，也终于成为他们中的一个。
